\originalchapter{Dirac 旋量与自旋求和}\label{ch:spinor}
\originalsection{一阶相对论方程}
能动量关系给 Klein--Gordon 方程 $(\Box+m^2)\phi=0$. 若希望得到对时间和空间都一阶的方程, 取 $(\ii\gamma^\mu\partial_\mu-m)\psi=0$, 平方后需要 $\{\gamma^\mu,\gamma^\nu\}=2g^{\mu\nu}$. Dirac 表象可用
\begin{equation}
 \gamma^0=\begin{pmatrix}I&0\\0&-I\end{pmatrix},\qquad
 \gamma^i=\begin{pmatrix}0&\sigma_i\\-\sigma_i&0\end{pmatrix}.
\end{equation}
$\sigma_i$ 是 Pauli 矩阵. 旋量的四个分量不是四种空间方向, 而是相对论自旋与粒子--反粒子结构的表示空间.

令 $\bar\psi=\psi^\dagger\gamma^0$, 电流 $j^\mu=\bar\psi\gamma^\mu\psi$ 满足连续性方程. $j^0=\psi^\dagger\psi$ 非负; $\bar\psi\psi$ 则是 Lorentz 标量而不是概率密度.

\originalsection{自由旋量与螺旋度}
平面波代入方程得 $(\slashed p-m)u=0$. 归一化 $u^\dagger u=2E$, 可取
\begin{equation}
 u(p,s)=\begin{pmatrix}\sqrt{E+m}\,\chi_s\\(\boldsymbol\sigma\cdot\mathbf p)\chi_s/\sqrt{E+m}\end{pmatrix}.
\end{equation}
反粒子旋量解 $(\slashed p+m)v=0$. 不同教材的 $v$ 自旋标签和整体相位可能不同; 比较显式电流之前必须先固定约定.

以 $h=\pm1$ 标记物理螺旋度的两倍, 方向 $\mathbf n=(\sin\theta\cos\varphi,\sin\theta\sin\varphi,\cos\theta)$ 可取
\begin{equation}
 \chi_+(\mathbf n)=\begin{pmatrix}\cos(\theta/2)\\\e^{\ii\varphi}\sin(\theta/2)\end{pmatrix},\quad
 \chi_-(\mathbf n)=\begin{pmatrix}-\e^{-\ii\varphi}\sin(\theta/2)\\\cos(\theta/2)\end{pmatrix}.
\end{equation}
无质量极限中 $u_h=\sqrt E(\chi_h,h\chi_h)^T$, 本稿反粒子约定取 $v_h=\sqrt E(-h\chi_{-h},\chi_{-h})^T$. 反粒子物理自旋标签与这里下两分量的 Pauli 标签相反. 这一约定使向量流只连接相反物理螺旋度的粒子反粒子对.

\originalsection{手征与弱相互作用}
$\gamma^5=\ii\gamma^0\gamma^1\gamma^2\gamma^3$, $P_{L,R}=(1\mp\gamma^5)/2$ 是手征投影. 对无质量粒子, 手征与螺旋度可对应; 对有质量粒子, boost 可改变运动方向与螺旋度, 但手征投影是场的代数结构. 把``弱作用只耦合左手征场''直接翻成``所有有质量粒子永远只允许一个自旋朝向''会失真.

Pion 为自旋零, 静止衰变要求两末态自旋合成零. 弱作用偏好左螺旋度中微子与右螺旋度反中微子, 带电轻子必须通过质量相关的手征混合匹配. 因而 $\pi\to e\nu$ 尽管相空间较大, 振幅仍有 $m_e$ 抑制, 宽度比约为
\begin{equation}
 \frac{\Gamma(\pi\to e\nu)}{\Gamma(\pi\to\mu\nu)}
 =\frac{m_e^2}{m_\mu^2}\frac{(m_\pi^2-m_e^2)^2}{(m_\pi^2-m_\mu^2)^2}.
\end{equation}

\originalsection{迹技巧}
未偏振反应用完备关系 $\sum_su_s\bar u_s=\slashed p+m$, $\sum_sv_s\bar v_s=\slashed p-m$. 例如
\begin{equation}
 \sum_{s_1,s_2}|\bar u_1\Gamma v_2|^2
 =\Tr[(\slashed p_1+m_1)\Gamma(\slashed p_2-m_2)\bar\Gamma],
 \quad\bar\Gamma=\gamma^0\Gamma^\dagger\gamma^0.
\end{equation}
它来自把旋量乘积的指标闭成矩阵迹, 不是假设每个螺旋度给相同振幅. 常用 $\Tr(\gamma^\mu\gamma^\nu)=4g^{\mu\nu}$ 与
$\Tr(\slashed a\gamma^\mu\slashed b\gamma^\nu)=4(a^\mu b^\nu+a^\nu b^\mu-g^{\mu\nu}a\cdot b)$. 奇数 gamma 的迹为零. 这些恒等式把四分量计算改写成不变量运算.

\originalsection{从二分量代数到轻子张量}
这里取$g=\operatorname{diag}(1,-1,-1,-1)$。Dirac分块中的矩阵须满足
\begin{equation}
\sigma_1=\begin{pmatrix}0&1\\1&0\end{pmatrix},\quad
\sigma_2=\begin{pmatrix}0&-\ii\\\ii&0\end{pmatrix},\quad
\sigma_3=\begin{pmatrix}1&0\\0&-1\end{pmatrix},\qquad
\sigma_i\sigma_j=\delta_{ij}I+\ii\epsilon_{ijk}\sigma_k .
\end{equation}
因此$(\boldsymbol\sigma\cdot\mathbf p)^2=\mathbf p^2I$，直接分块相乘可核实$\{\gamma^\mu,\gamma^\nu\}=2g^{\mu\nu}$。斜杠符号定义为$\slashed a=\gamma^\mu a_\mu=\gamma^0a^0-\boldsymbol\gamma\cdot\mathbf a$。后文的协变导数与场强则分别是$\ii D_\mu=\ii\partial_\mu-qA_\mu$和$F_{\mu\nu}=\partial_\mu A_\nu-\partial_\nu A_\mu$；$\slashed D=\gamma^\mu D_\mu$，$F$不是另一种旋量。$q$包括电荷符号，换成电子的$-e$约定时须一同改变顶点符号。

写$u=(\phi,\eta)^T$，$(\slashed p-m)u=0$给$(E-m)\phi=(\boldsymbol\sigma\cdot\mathbf p)\eta$、$(E+m)\eta=(\boldsymbol\sigma\cdot\mathbf p)\phi$。取$\phi=\sqrt{E+m}\chi_s$即得上节$u$；取相反能量方程的正能反粒子标记则得
\begin{equation}
v(p,s)=\begin{pmatrix}(\boldsymbol\sigma\cdot\mathbf p)\xi_s/\sqrt{E+m}\\\sqrt{E+m}\xi_s\end{pmatrix},
\qquad \bar u_su_{s'}=2m\delta_{ss'},\quad\bar v_sv_{s'}=-2m\delta_{ss'} .
\end{equation}
这里$\xi_s$只是Pauli基矢；上节$v_h$用物理螺旋度重新标号，相位约定可不同。把$\sum_s\chi_s\chi_s^\dagger=\sum_s\xi_s\xi_s^\dagger=I$代入四个$2\times2$分块外积，并用$(\boldsymbol\sigma\cdot\mathbf p)^2=\mathbf p^2$，得到
\begin{equation}
\sum_su_s\bar u_s=\slashed p+m,\qquad\sum_sv_s\bar v_s=\slashed p-m.
\end{equation}
这给出上节完备关系的来历，也说明$\bar uu$和$u^\dagger u$为何分别是$2m$与$2E$。

再把迹规则实际用于下一章的$e^-(p_1)e^+(p_2)\to\mu^-(k_1)\mu^+(k_2)$。质量可忽略时定义
\begin{equation}
L^{\mu\nu}=\Tr[\slashed p_2\gamma^\mu\slashed p_1\gamma^\nu],\qquad
H_{\mu\nu}=\Tr[\slashed k_1\gamma_\mu\slashed k_2\gamma_\nu].
\end{equation}
每个初态自旋各有两种，末态已求和，故平均因子是$1/4$。展开并用$p_1\cdot k_1=p_2\cdot k_2=-t/2$、$p_1\cdot k_2=p_2\cdot k_1=-u/2$，得到$L^{\mu\nu}H_{\mu\nu}=8(t^2+u^2)$及
$(e^4/4s^2)L^{\mu\nu}H_{\mu\nu}=2e^4(t^2+u^2)/s^2$。这个收缩把下一章角分布中的每个自旋平均因子公开给读者核查。
